% Generated draft: requires mathematical and source review.
\documentclass[11pt]{article}
\usepackage[margin=25mm]{geometry}
\usepackage{fontspec}
\setmainfont{Latin Modern Roman}
\usepackage{amsmath,amssymb,mathtools,unicode-math}
\setmathfont{Latin Modern Math}
\usepackage{enumitem,needspace,xurl,hyperref}
\setlength{\parindent}{0pt}
\setlength{\parskip}{5pt}
\setlength{\emergencystretch}{3em}
\newcommand{\N}{\mathbb{N}}
\newcommand{\Z}{\mathbb{Z}}
\newcommand{\Q}{\mathbb{Q}}
\newcommand{\R}{\mathbb{R}}
\newcommand{\C}{\mathbb{C}}
\newcommand{\F}{\mathbb{F}}
\newcommand{\A}{\mathbb{A}}
\newcommand{\PP}{\mathbb{P}}
\newcommand{\E}{\mathbb{E}}
\newcommand{\eps}{\varepsilon}
\newcommand{\dd}{\,\mathrm{d}}
\begin{document}
\section*{UnsolvedMath ID: 30001999}
\textbf{Source problem code:} catalogue entry
\subsection{Definitions and mathematical statement}
The supplied record names a critical-level simple module \(L(\lambda)\), but does not specify the affine Lie algebra, its normalization, the category of modules, or the allowed highest weights \(\lambda\). The term \emph{critical level} refers to a distinguished value of the central action; its numerical value cannot be fixed here without the missing algebra and normalization conventions.

For a specified Cartan subalgebra \(\mathfrak h\), a weight module has a decomposition
\[
L(\lambda)=\bigoplus_{\mu\in\mathfrak h^*}L(\lambda)_\mu,
\qquad
L(\lambda)_\mu=\{v\in L(\lambda):hv=\mu(h)v\text{ for every }h\in\mathfrak h\}.
\]
If these weight spaces are finite-dimensional, its formal character is
\[
\operatorname{ch}L(\lambda)=\sum_{\mu\in\mathfrak h^*}\dim L(\lambda)_\mu\,e^\mu,
\]
where \(e^\mu\) denotes a formal weight symbol. These are conditional definitions: the record does not establish the relevant weight-space finiteness or specify the formal completion in which a character identity is intended.

\textbf{Source-stated problem.} Prove the Feigin--Frenkel--Lusztig character formula expressing the character of a critical-level simple module \(L(\lambda)\) in terms of periodic polynomials.

\textbf{Unresolved formulation.} The supplied record gives neither the character identity nor the definition, indexing, normalization, or evaluation of the periodic polynomials. It also gives no precise range of \(\lambda\). Consequently, the requested identity and its quantifiers cannot be reconstructed uniquely from this record. The supplied background additionally describes an unfinished conjectural program relating sheaves in an affine flag setting to positive-level category \(\mathcal O\), without specifying its connection to the critical-level assertion. A complete mathematical formulation requires these missing specifications.
\subsection{Short English statement}
The recorded problem asks for a character formula for critical-level simple modules using periodic polynomials. The algebra, permitted weights, polynomial conventions, and actual formula are missing, so this remains a statement-only draft requiring source clarification and human mathematical review.
\subsection{Sources}
\begin{itemize}
\item Local UnsolvedMath registry, permanent ID \(30001999\), source problem code \texttt{catalogue entry}, title ``Feigin--Frenkel--Lusztig Character Formula.'' The supplied record belongs to the Oberwolfach Reports Open Problems collection and gives the proposed year as 2012. No individual proposer is supplied; the names in the formula title do not establish proposer attribution.
\item \emph{Enveloping Algebras and Geometric Representation Theory} (2012), Oberwolfach Reports, DOI 10.4171/owr/2012/13, pp. 733--809. Supplied source link and literature-source entry ``Oberwolfach Report 13 2012'': \url{https://doi.org/10.4171/owr/2012/13}. The record describes these reports as documenting workshops and their open-problem discussions. No item-specific page or passage is supplied.
\item Supplied literature assessment, checked on \texttt{2026-08-21}: status ``open.'' The assessment describes an unfinished conjectural program linking sheaves in the affine flag setting to positive-level category \(\mathcal O\), and reports that the available literature does not establish the stated general critical-representation consequences. This is supplied dated curation, not an independently verified literature assessment; the record explicitly requests independent verification of time-sensitive status and cited claims.
\item The registry labels the statement \texttt{corrected\_verified} and reports a manual semantic, typographic, and LaTeX review confirming a meaning-preserving cleanup. This supplied editorial label does not resolve the mathematical specifications missing from the record.
\end{itemize}
\end{document}
